\documentclass[11pt]{article}

\usepackage[margin=1in]{geometry}
\usepackage{booktabs}
\usepackage{hyperref}
\usepackage{amsmath}
\usepackage{amssymb}
\usepackage{array}
\usepackage{xcolor}

\hypersetup{
  colorlinks=true,
  linkcolor=blue!45!black,
  urlcolor=blue!45!black,
  citecolor=blue!45!black
}

\newcommand{\TzEL}{TzEL}
\newcommand{\Hash}{\mathsf{H}}
\newcommand{\cm}{\mathsf{cm}}
\newcommand{\nf}{\mathsf{nf}}
\newcommand{\rcm}{\mathsf{rcm}}

\title{\TzEL: Tezos Encrypted Ledger\\
\large Post-Quantum Shielded Transactions on Tezos Smart Rollups}
\author{TzEL contributors}
\date{}

\begin{document}
\maketitle

\begin{abstract}
\TzEL{}, the Tezos Encrypted Ledger, is a shielded transaction protocol for
Tezos smart rollups. It is designed for a setting in which shielded transaction
data is public and durable, so encrypted note payloads may be archived long
before anyone knows whether the cryptography protecting them will survive
future attacks. \TzEL{} adapts the note/commitment/nullifier architecture of
shielded payments to that threat by making the shielded path post-quantum end
to end: incoming notes are delivered with ML-KEM, spends are authorized with
one-time hash-based signatures inside an XMSS-style tree, and validity is
proved with recursive STARKs. The protocol also separates spend authority from
proof generation, allowing wallets to authorize exact transaction effects
locally while outsourcing proof generation to an untrusted prover. Shielded
transactions can move value from bridge deposits into notes, transfer value
between notes, and release value back to public recipients. The result is a
rollup-based shielded ledger whose archived note ciphertexts, spend
authorizations, and proof system do not depend on classical discrete-log
assumptions.
\end{abstract}

\section{Introduction}

Shielded payment systems publish encrypted transaction data on a public ledger.
Those ciphertexts can be copied and stored indefinitely. If the cryptography
protecting them is later broken, old notes, values, and memos can be decrypted
retroactively. For a private payment system, that is the natural form of
harvest-now-decrypt-later.

That observation creates two distinct questions. The first is \emph{when} to
migrate: if note ciphertexts can be harvested now and decrypted later, then the
privacy layer matters before a quantum-capable adversary arrives. The second is
\emph{how far} to migrate: it is possible to harden only the memo layer and
still leave other parts of the shielded path on assumptions that are not meant
to survive quantum attacks. \TzEL{} takes the broader route. It aims to make
the shielded transaction path itself post-quantum.

A second requirement is practical rather than cryptographic. STARK proving is
heavy enough that wallets may want to outsource proof generation. That should
not require outsourcing spend authority. \TzEL{} therefore separates the right
to authorize a spend from the work of producing the proof that validates it.

\section{Threat Model and Design Goals}

\subsection{Threat model}

\TzEL{} is designed for the following setting:

\begin{itemize}
  \item Shielded transaction data, including encrypted note payloads, is public
        and can be archived indefinitely.
  \item Proof generation may be delegated to an untrusted prover.
  \item Users may want watch-only services that can help find incoming notes or
        account for sent outputs without gaining spend authority.
  \item The rollup enforces public state-transition rules, while the proof
        system enforces private consistency of notes, nullifiers, and spends.
\end{itemize}

Like other practical shielded ledgers, \TzEL{} does not make all observable
structure disappear. Transaction type, timing, ordering, number of spent
inputs, and some output-shape information remain visible at the public layer.
The protocol separates receiver-side recovery from sender-side recovery.
Recipients use incoming viewing material to detect and decrypt notes addressed
to them. Senders can separately export outgoing viewing material that recovers
the metadata needed to recognize notes they created, without granting spend
authority and without adding another public key to payment addresses.

\subsection{Design goals}

The main design goals are:

\begin{itemize}
  \item Preserve the core shielded-ledger structure of notes, commitments,
        nullifiers, and historical Merkle roots.
  \item Replace non-post-quantum assumptions in the shielded path with
        primitives chosen for long-lived public ciphertexts and durable note
        archives.
  \item Keep spend authority in the wallet even when proof generation is
        delegated.
  \item Support reduced-capability exports such as detection-only and
        incoming-view monitoring and sender-side outgoing accounting.
  \item Fit the execution and data model of Tezos smart rollups and the Tezos
        Data Availability Layer (DAL).
\end{itemize}

\section{Protocol Overview}

A \emph{note} is the protocol's private value object. On chain it appears as a
commitment together with encrypted note data. Off chain, the recipient uses the
viewing path and local address metadata to recover the note plaintext. To spend
a note, the wallet reveals a \emph{nullifier}: a one-use serial derived from
the note and the note owner's secret spending material. The rollup rejects any
nullifier that has already appeared.

The rollup state therefore contains:

\begin{itemize}
  \item an append-only Merkle tree of note commitments;
  \item a global set of spent nullifiers;
  \item a historical set of Merkle roots usable as spend anchors;
  \item a set of public bridge deposit receipts keyed by secret-bound deposit
        identifiers, used only to move value into the shielded pool;
  \item verifier configuration, public-account bindings, and fee-policy state.
\end{itemize}

A \emph{shielded transaction} is any transaction that creates or destroys notes.
\TzEL{} has three of them:

\begin{center}
\begin{tabular}{>{\raggedright\arraybackslash}p{0.17\linewidth} >{\raggedright\arraybackslash}p{0.28\linewidth} >{\raggedright\arraybackslash}p{0.42\linewidth}}
\toprule
Transaction & Consumes & Creates \\
\midrule
Shield & a bridge deposit receipt keyed by a secret-bound deposit identifier & a user note and a DAL-producer fee note \\
Transfer & 1--7 shielded input notes & a recipient note, a change note, and a DAL-producer fee note \\
Unshield & 1--7 shielded input notes & a public credit, an optional change note, and a DAL-producer fee note \\
\bottomrule
\end{tabular}
\end{center}

Every shielded transaction pays two distinct resource prices. The first is a
burned rollup fee \(fee\), which prices rollup execution and is recreated
nowhere. In the current reference rollup it has a floor of 100000 mutez: the
first two private transactions in an inbox level pay the floor, each later
private transaction doubles the required fee, the doubling is capped after six
steps, and the fee resets when the inbox level advances. The second is a
strictly positive private producer fee \(producer\_fee\), created as an
ordinary shielded note intended for the DAL slot producer. The rollup proves
that such a note exists and has positive value; the producer enforces that it is
payable to them as an inclusion policy before publishing. The burned fee and the
producer fee price different resources and are therefore modeled separately.

\section{Cryptographic Design}

\subsection{Post-quantum profile}

\begin{center}
\begin{tabular}{>{\raggedright\arraybackslash}p{0.28\linewidth} >{\raggedright\arraybackslash}p{0.31\linewidth} >{\raggedright\arraybackslash}p{0.27\linewidth}}
\toprule
Role & Instantiation & Purpose \\
\midrule
Commitments and nullifiers & BLAKE2s-derived field hashes & bind notes and prevent double spends \\
Incoming-note delivery & ML-KEM-768 + ChaCha20-Poly1305 & encrypt note payloads and memos \\
Detection & ML-KEM-768 clue + short tag & watch-only candidate filtering \\
Outgoing recovery & BLAKE2s-derived key + ChaCha20-Poly1305 & sender-side accounting for created outputs \\
Spend authorization & hash-based one-time signatures in an XMSS-style tree & authorize each spent note once \\
Zero knowledge & Cairo AIR + recursive STARKs & prove validity without exposing witness data \\
\bottomrule
\end{tabular}
\end{center}

Where the protocol can be expressed with hash-based machinery, it does so:
commitments, nullifiers, Merkle nodes, owner binding, transaction sighashes,
and one-time-signature chains are all built from hashes. That leaves one place
where hashing alone is not enough. A sender must take a public payment address
and, without interaction, produce note ciphertext that only the recipient can
open. That is a public-key delivery problem, so \TzEL{} uses ML-KEM for the
incoming-note path.

This is why the system is not purely hash-based. If there were a practical,
standardized hash-based KEM with the right ergonomics for per-address note
delivery and watch-only detection, the protocol could have used it. In the
current ecosystem, the pragmatic post-quantum choice for that role is ML-KEM.
Everywhere else, the design stays with hashes.

That still leaves a design choice. One could adopt ML-KEM for note encryption
and stop there, while keeping spend authorization or proof verification on
classical assumptions. \TzEL{} takes the stronger route: the shielded path as a
whole is meant to survive without elliptic-curve or pairing-based cryptography.

\subsection{Proof architecture}

The private transaction logic is first expressed as Cairo AIR. That first proof
establishes the arithmetic and Merkle constraints for shield, transfer, or
unshield. A second proving layer then re-proves that Cairo execution using
Stwo, yielding the recursive proof object that is posted to the rollup.

This two-layer structure serves two purposes. First, it keeps the proving stack
entirely in STARK territory. Second, it yields a proof that is small and
verifier-friendly enough for rollup use while still carrying zero-knowledge
blinding. On the current reference stack, recursive proofs are around 300 KB;
single-level proving exists as a development aid, not as the intended privacy
profile.

\section{Keys, Addresses, and Capabilities}

The key hierarchy is organized around capability separation. \TzEL{}
distinguishes the following reduced-authority capabilities:

\begin{center}
\begin{tabular}{>{\raggedright\arraybackslash}p{0.18\linewidth} >{\raggedright\arraybackslash}p{0.60\linewidth}}
\toprule
Capability & What it can do \\
\midrule
Detection & flag candidate incoming notes, with false positives, without reading note contents \\
Incoming view & decrypt and validate incoming notes and memos, but not spend them \\
Outgoing view & recover sender-created output metadata and values, but not spend them \\
Full view & incoming view plus spent-state tracking for notes whose address metadata is known \\
Spend & full view plus transaction authorization \\
\bottomrule
\end{tabular}
\end{center}

Detection, incoming view, full view, and spend form a nested receiver-side
chain. Outgoing view is orthogonal: it tracks outputs created by the sender but
does not detect arbitrary incoming notes, compute nullifiers, or authorize
spends.

That separation is reflected in the wallet derivation tree. From a master
secret, the wallet derives independent spending, incoming, and outgoing
branches:

\[
\begin{aligned}
spend\_seed &= \Hash(\mathrm{TAG\_SPEND}, master\_sk), \\
incoming\_seed &= \Hash(\mathrm{TAG\_INCOMING}, master\_sk), \\
outgoing\_seed &= \Hash(\mathrm{TAG\_OUTGOING}, master\_sk).
\end{aligned}
\]

The spending branch contains the secrets needed to derive nullifiers and to
authorize spends. The incoming branch contains the material needed to receive,
detect, and view notes. This makes reduced-capability export natural: a user
can hand detection or viewing material to another process without handing over
spend authority.

Incoming and full-view wallets still need the local address records containing
\((d_j, auth\_root_j, pub\_seed_j, nk\_tag_j)\) to validate recovered notes.
The incoming and nullifier branches alone do not reconstruct that public address
metadata.

The outgoing branch is sender-scoped. It is not part of payment addresses and
does not require an additional ML-KEM public key on chain. Instead, when a wallet
creates an output, it encrypts a compact recovery record under a symmetric key
derived from \(outgoing\_seed\) and the output commitment \(\cm\). This lets the
sender later recognize and account for notes it created while preserving the
recipient-facing address format.

For address index \(j\), the wallet derives:

\begin{itemize}
  \item \(d_j\), the address diversifier;
  \item \(nk\_spend_j\), the per-address nullifier secret;
  \item \(nk\_tag_j = \Hash_{\mathrm{nktg}}(nk\_spend_j)\), a public binding tag
        that lets a sender bind a note to the correct nullifier lineage without
        learning the nullifier secret;
  \item \((auth\_root_j, pub\_seed_j)\), the public root and seed of that
        address's one-time authorization tree;
  \item \((ek^v_j, dk^v_j)\), the ML-KEM viewing keypair;
  \item \((ek^d_j, dk^d_j)\), the ML-KEM detection keypair.
\end{itemize}

The payment address is therefore

\[
address_j = (d_j, auth\_root_j, pub\_seed_j, nk\_tag_j, ek^v_j, ek^d_j).
\]

It contains enough public information for a sender to create a spendable note
for that address, but not enough information to derive nullifiers or spend
authorizations.

These addresses are larger than classical shielded addresses because they carry
multiple post-quantum public components. In practice, such addresses are better
handled as structured payment requests, QR encodings, or directory lookups than
as strings to copy by hand.

\section{Notes, Commitments, and Nullifiers}

For note value \(v\), per-note seed \(rseed\), and the recipient's address
metadata, the protocol derives commitment randomness and the note commitment as

\[
\rcm = \Hash(\Hash(\mathrm{TAG\_RCM}), rseed),
\]

\[
owner\_tag_j = \Hash_{\mathrm{owner}}(auth\_root_j, pub\_seed_j, nk\_tag_j),
\]

\[
\cm = \Hash_{\mathrm{commit}}(d_j, v, \rcm, owner\_tag_j).
\]

The owner tag is what fuses the two sides of future spending authority into the
note commitment. It binds the note both to the authorization tree and to the
nullifier derivation path. Without that binding, a commitment could be paired
with unrelated nullifier material. The dependency chain is

\[
nk\_spend_j \rightarrow nk\_tag_j \rightarrow owner\_tag_j \rightarrow \cm.
\]

Nullifiers are position-dependent:

\[
\nf = \Hash_{\mathrm{nf}}(nk\_spend_j, \Hash_{\mathrm{nf}}(\cm, pos)).
\]

Including the Merkle leaf position prevents two otherwise identical commitments
inserted at different leaves from collapsing to the same nullifier. That is the
same basic issue sometimes described as a faerie-gold problem: identical-looking
notes at different positions must still spend distinctly.

\section{Spend Authorization and Delegated Proving}

\subsection{Why one-time authorization is needed}

If proof generation is delegated, the wallet needs a way to authorize the exact
public effects of a transaction without handing the prover a reusable spending
key. \TzEL{} solves that with an XMSS-style authorization tree of one-time
hash-signature keys. Each spent input consumes one unused leaf. The wallet
signs the transaction locally; the prover only proves, inside the STARK, that
the signature is valid under the address's authorization root.

This mechanism is still useful when proving is done locally. It keeps spend
authorization inside the proof and avoids publishing public keys or signatures
on chain.

\subsection{Authorization tree and in-circuit verification}

Each address owns a depth-16 authorization tree, giving \(2^{16}\) one-time
authorization leaves. The per-leaf signature is WOTS-like, with the current
parameterization using base \(w = 4\), 133 chains, and chain length \(w - 1\).

For each spent input, the wallet signs a transaction-specific sighash with the
corresponding one-time secret. Inside the STARK, the circuit

\begin{enumerate}
  \item recomputes the sighash from the public outputs;
  \item decomposes it into base-4 WOTS digits;
  \item hashes the signature chains forward to recover the corresponding public
        key endpoints;
  \item compresses those endpoints through the XMSS-style L-tree;
  \item proves membership of the recovered leaf under \(auth\_root_j\).
\end{enumerate}

No authorization leaf, one-time public key, or one-time signature appears in the
public transaction outputs. The proof itself is the public authorization object.

That matters for privacy as well as correctness. Publishing authorization leaves
or signatures would create another public handle linking spends to address
material. Instead, the public layer sees only the commitments, nullifiers,
note-data hashes, and other public effects that the rollup actually needs.

\subsection{What the prover can and cannot do}

The wallet, not the prover, chooses the transaction effects. It selects the
inputs, computes the output commitments, derives the note-data hashes, forms
the public statement, and signs the sighash locally. These hashes bind the
posted recipient-delivery ciphertexts and outgoing-recovery ciphertexts as
bytes; the circuit does not prove that those ciphertexts decrypt to the
commitment preimages. The prover receives only the witness needed to prove that
statement.

Because spend authorization is checked inside the circuit over the public
outputs, an untrusted prover cannot change recipients, fees, nullifiers,
commitments, withdrawal targets, or note-data hashes without invalidating the
proof. This is the core separation between delegated proving and delegated
spending.

\subsection{Replay binding}

For transfer and unshield, the signature covers a folded transaction sighash
that includes

\begin{itemize}
  \item a circuit-type tag;
  \item the deployment authorization domain \(auth\_domain\);
  \item the input root;
  \item all nullifiers;
  \item public values such as fees, withdrawal amount, and recipient id;
  \item output commitments and note-data hashes, including the hashes of
        recipient-delivery and outgoing-recovery ciphertexts.
\end{itemize}

The type tag prevents cross-circuit replay. The authorization domain prevents
replay across mirrored deployments, forks, or verifier migrations that might
otherwise share root history. The current base protocol does not yet include an
expiry or per-transaction nonce in that authorization, so a delegated prover
can withhold a completed authorization until one of its nullifiers is consumed
elsewhere.

\section{Shielded Transaction Types}

\subsection{Shield}

Shield takes value that has already crossed from L1 into the rollup and moves it
into the shielded pool. The public rollup object being consumed is not a normal
transparent account balance. It is a bridge deposit receipt keyed by a
secret-bound deposit identifier:

\[
deposit\_id = \Hash(\mathrm{TAG\_DEPOSIT}, deposit\_secret).
\]

The bridge credits a public rollup balance key derived from that
\(deposit\_id\). To shield the funds, the witness proves knowledge of
\(deposit\_secret\) and therefore of the exact deposit receipt being consumed.

The public outputs of shield are

\[
[v_{pub}, fee, producer\_fee, \cm_{user}, \cm_{producer}, deposit\_id,
 mh_{user}, mh_{producer}].
\]

The circuit proves the user commitment, the producer-fee commitment, that
\(deposit\_id\) equals \(\Hash(\mathrm{TAG\_DEPOSIT}, deposit\_secret)\), and
the condition \(producer\_fee > 0\). The rollup then checks proof validity,
deposit binding, note-data hashes, and the required burned fee.

The consumed public balance is

\[
v_{pub} + fee + producer\_fee.
\]

That accounting is deliberate. \(v_{pub}\) becomes the user's shielded note,
\(producer\_fee\) becomes the DAL-producer fee note, and \(fee\) is burned. The
deposit receipt itself is not directly withdrawable; otherwise publication of
\(deposit\_id\) would be enough to race the shield and defeat the ownership
binding.

\subsection{Transfer}

Transfer consumes \(N\) shielded notes, where \(1 \leq N \leq 7\), and creates
three new shielded notes: recipient, change, and DAL-producer fee. Its public
outputs are

\[
[auth\_domain, root, \nf_0, \ldots, \nf_{N-1}, fee,
 \cm_1, \cm_2, \cm_3, mh_1, mh_2, mh_3].
\]

For each input, the circuit proves commitment reconstruction, Merkle membership,
nullifier computation, and one-time authorization under the corresponding
authorization root. For the outputs, it proves all three commitments and
enforces

\[
\sum_{i=0}^{N-1} v_i = v_1 + v_2 + v_3 + fee,
\]

with \(v_3 > 0\) for the producer-fee note.

The rollup state transition checks the public nullifier list before applying the
transfer. After validating and stripping the bootloader's Cairo array length
prefix, the verified public-output vector must have exactly the length implied by
\(N\), the published nullifiers must be pairwise distinct, and none of them may
already appear in the global nullifier set.

\subsection{Unshield}

Unshield consumes \(N\) shielded notes, releases public value to a recipient,
and creates a producer-fee note plus an optional shielded change note. Its
public outputs are

\[
[auth\_domain, root, \nf_0, \ldots, \nf_{N-1}, v_{pub}, fee,
 recipient\_id, \cm_{change}, mh_{change}, \cm_{fee}, mh_{fee}].
\]

The circuit proves the same input validity and authorization conditions as
transfer, then enforces

\[
\sum_{i=0}^{N-1} v_i = v_{pub} + v_{change} + v_{fee} + fee.
\]

If there is no change note, the change commitment, change note-data hash, and all
corresponding witness values are constrained to zero. That removes a source of
prover malleability: the prover cannot hide nonzero private change behind a
zero public commitment.

The rollup applies the same public nullifier checks as transfer: after validating
and stripping the bootloader's Cairo array length prefix, the verified
public-output vector length must match \(N\), the published nullifiers must be
pairwise distinct, and every published nullifier must be fresh with respect to
the global nullifier set.

\section{Encrypted Note Delivery and Detection}

Each output note carries

\begin{itemize}
  \item the commitment \(\cm\);
  \item a detection ciphertext \(ct_d\);
  \item a short detection tag;
  \item a viewing ciphertext \(ct_v\);
  \item a derived AEAD nonce;
  \item encrypted note plaintext \((v \parallel rseed \parallel memo)\);
  \item an outgoing recovery ciphertext \(ct_o\).
\end{itemize}

For recipient address \((ek^d_j, ek^v_j)\), the sender creates two ML-KEM
encapsulations. The viewing path is the strong one: encapsulate to \(ek^v_j\),
derive a symmetric key from the shared secret, and encrypt
\((v \parallel rseed \parallel memo)\) with ChaCha20-Poly1305. The detection
path is weaker on purpose. It gives a watch-only service just enough material
to flag candidate incoming notes, while leaving note contents encrypted and
spend authority untouched.

This is the operational role of detection. Without it, a wallet or service would
have to attempt full incoming-note recovery against every output on the chain.
With detection, a lightweight indexer can scan the chain, report likely
matches by address index, and leave the wallet to do full note recovery only on
that narrowed candidate set.

Detection leaks very little compared with viewing. A detector can test whether a
note is a candidate for one of the monitored addresses. It cannot decrypt note
contents, derive nullifiers, authorize spends, or directly infer spent status.
False positives are built in by design.

The outgoing recovery path is separate from recipient delivery. For each output,
the sender encrypts

\[
(role \parallel v \parallel rseed \parallel d_j \parallel auth\_root_j
 \parallel pub\_seed_j \parallel nk\_tag_j)
\]

under a key derived from \(outgoing\_seed\) and \(\cm\). The plaintext omits
spending secrets. A wallet using outgoing viewing material accepts a recovered
record only after recomputing

\[
\cm_{expected} =
\Hash_{\mathrm{commit}}(d_j, v, \Hash(\Hash(\mathrm{TAG\_RCM}), rseed),
\Hash_{\mathrm{owner}}(auth\_root_j, pub\_seed_j, nk\_tag_j))
\]

and checking \(\cm_{expected} = \cm\). Thus outgoing viewing is an accounting
and recovery feature for the sender, not a new authorization path.

Detection is only a filtering aid, not a correctness condition. A wallet accepts
an incoming note only after it

\begin{enumerate}
  \item decrypts the note plaintext;
  \item selects the matching local address record containing
        \((d_j, auth\_root_j, pub\_seed_j, nk\_tag_j)\);
  \item recomputes \(\rcm = \Hash(\Hash(\mathrm{TAG\_RCM}), rseed)\);
  \item recomputes
        \(owner\_tag_j = \Hash_{\mathrm{owner}}(auth\_root_j, pub\_seed_j, nk\_tag_j)\);
  \item recomputes
        \(\cm_{expected} = \Hash_{\mathrm{commit}}(d_j, v, \rcm, owner\_tag_j)\);
  \item accepts the note only if \(\cm_{expected}\) equals the posted
        commitment.
\end{enumerate}

These acceptance rules matter because the proof binds the posted ciphertext
bytes through note-data hashes, but does not prove that the viewing or outgoing
ciphertexts decrypt to the same commitment preimage used inside the circuit.

\section{Rollup Verification and DAL}

The STARK proves private consistency. The rollup still has to enforce public
state-transition rules that are not handled inside the proof. These include:

\begin{itemize}
  \item proof verification for the intended executable;
  \item \(auth\_domain\) matching the configured deployment domain for
        spending transactions;
  \item the input root being a recognized historical root;
  \item exact agreement between the verified public-output length and the
        number of published nullifiers;
  \item pairwise distinctness and global freshness of every published
        nullifier;
  \item equality between posted note commitments and the commitments in the
        proof outputs;
  \item equality between posted note data and the note-data hashes in the
        proof outputs;
  \item correct binding of shield to the exact secret-bound deposit receipt
        being consumed;
  \item correct binding of unshield to the intended public recipient encoding;
  \item enough remaining commitment-tree capacity for every output note before
        any balance, nullifier, or outbox mutation is applied;
  \item satisfaction of the current burned-fee policy.
\end{itemize}

These checks must be implemented and should be treated as part of the protocol.
Without historical-root validation, a prover could spend against a fabricated local
tree. Without executable binding, a verifier could accept the wrong Cairo task.
Without note-data binding, a relayer could swap ciphertext fields after the
proof was produced.

Verifier parameters are not supplied by the transaction. The rollup infers the
expected circuit from the operation type, checks the corresponding configured
executable hash, and derives the verifier inputs from the proof output
preimage.

The large proof and note payloads are why the Tezos DAL is structurally part of
the design. The DAL is a companion data-availability network for rollups: it
lets rollups publish large blobs without forcing them into L1 block bandwidth.
That matters here because recursive proofs are roughly 300 KB and each encrypted
output note is roughly 3 KB. The intended path is therefore to publish the
heavy payload through DAL, inject only a small pointer into the rollup inbox,
and let the kernel resolve that pointer before applying the state transition.

\TzEL{}'s post-quantum claims concern the shielded transaction layer. Tezos L1
consensus, bridge mechanics, and DAL attestation are separate systems with
their own security assumptions. Even before those layers reach a fully
post-quantum profile, protecting note ciphertexts, spend authorizations, and
proof objects at the shielded layer still blocks the harvest-now-decrypt-later
attack against archived shielded data.

\section{Security Properties and Limitations}

\subsection{What the design gives you}

\begin{itemize}
  \item \textbf{Balance conservation.} Values are range-checked and the value
        equations for shield, transfer, and unshield are enforced inside the
        circuit.
  \item \textbf{Double-spend resistance.} Nullifiers are unique per
        note-position pair. The circuit binds each public nullifier to the
        corresponding private input witness, and the rollup enforces that the
        public nullifier list has no duplicates and is disjoint from the global
        spent-nullifier set before applying the transaction.
  \item \textbf{Commitment binding.} The chain
        \(nk\_spend \rightarrow nk\_tag \rightarrow owner\_tag \rightarrow \cm\)
        ties each note to the nullifier lineage that will later be used to
        spend it.
  \item \textbf{Local spend control.} A valid spend proof implies both note
        ownership and a valid one-time authorization over the exact public
        effects of the transaction.
  \item \textbf{No public authorization artifacts.} Authorization leaves,
        one-time public keys, and one-time signatures do not appear in the
        public transaction outputs.
  \item \textbf{Post-quantum shielded path.} Incoming-note delivery, spend
        authorization, and zero-knowledge proving do not depend on classical
        discrete-log assumptions.
\end{itemize}

\subsection{What still remains visible}

\begin{itemize}
  \item Transaction type, ordering, and timing are public.
  \item The number of published nullifiers reveals the number of spent inputs.
  \item Observers can tell whether an unshield carries a change note.
  \item A delegated prover sees per-address witness material such as
        \(nk\_spend_j\) and \(auth\_root_j\). Combined with public commitments,
        positions, and the global nullifier set, that is enough to infer
        spent-state for notes associated with that address.
\end{itemize}

Users who want delegated proving with weaker witness exposure can treat TEE-based
provers or equivalent isolation as an operational mitigation, but that is not a
protocol guarantee.

\subsection{Current caveats}

\begin{itemize}
  \item \textbf{Detection is honest-sender.} A malicious sender can post bogus
        detection data for a note. That does not let them steal funds or read
        wallet state; it simply forces the recipient or its watch service to
        fall back to broader viewing-key scanning. Detection should therefore
        be understood as a performance optimization, not as the source of note
        correctness.
  \item \textbf{Ciphertext correctness is not proved in-circuit.} Wallets must
        recompute commitments before treating incoming or outgoing recovered
        records as valid.
  \item \textbf{Address fields are not self-authenticating to the sender.}
        Shield and transfer outputs can be formed with malformed public address
        metadata, producing unspendable notes. This is sender self-griefing,
        not theft.
  \item \textbf{One-time keys must still be treated as one-time.} Reuse does
        not create a public on-chain linking issue by itself because signatures
        are not published. But once the same external prover sees two
        authorizations under the same one-time key, it may be able to forge
        further authorizations. Wallets should therefore allocate one leaf per
        spend and persist that state durably.
  \item \textbf{Shield is not yet domain-bound the way transfer and unshield
        are.} Because shield carries \(deposit\_id\) but no \(auth\_domain\),
        mirrored deployments with matching funded deposit identifiers still
        require external deployment separation.
\end{itemize}

\section{Size and Data Path}

Each encrypted output note is roughly 3487 bytes on chain:

\begin{itemize}
  \item 32 bytes for the commitment;
  \item 1088 bytes for the detection ciphertext;
  \item 2 bytes for the detection tag;
  \item 1088 bytes for the viewing ciphertext;
  \item 12 bytes for the AEAD nonce;
  \item 1080 bytes for the encrypted \((v, rseed, memo)\) payload;
  \item 185 bytes for the outgoing recovery ciphertext.
\end{itemize}

Typical transaction sizes are therefore dominated by the recursive proof rather
than by authorization artifacts:

\begin{center}
\begin{tabular}{>{\raggedright\arraybackslash}p{0.20\linewidth} >{\raggedright\arraybackslash}p{0.22\linewidth} >{\raggedright\arraybackslash}p{0.36\linewidth}}
\toprule
Transaction & Output note data & Approximate total \\
\midrule
Shield & 2 notes & about 302 KB \\
Transfer & 3 notes & about 305 KB plus 32 bytes per input nullifier \\
Unshield & 1--2 notes & about 299--302 KB plus 32 bytes per input nullifier \\
\bottomrule
\end{tabular}
\end{center}

Those numbers are why DAL is not an optional add-on for this design. A rollup
that wants shielded transactions with recursive proofs needs a high-bandwidth
data path, not only a small inbox.

\section{Conclusion}

\TzEL{} preserves the core shielded-ledger logic of notes, commitments,
nullifiers, and Merkle-root-anchored spends, but retools the privacy layer for
a post-quantum setting and a rollup execution model. Hash-based constructions
cover commitments, nullifiers, and one-time spend authorization wherever they
can. ML-KEM is used where the protocol genuinely needs public-key delivery of
incoming notes and watch-only detection. Recursive STARKs replace pairing-based
proof systems while keeping spend authorization inside the proof itself.

The result is a shielded transaction protocol for Tezos smart rollups that aims
to protect the data most exposed to harvest-now-decrypt-later pressure: durable
note ciphertexts, spend authorizations, and the proof path that validates
private state transitions.
\end{document}
